%! Polynomial Functions and End Behavior — Unit 2, Lesson 2.5 — Lesson Plan
\documentclass[10pt]{article}
\usepackage{precalculus-article}
\usepackage{precalculus-boxes}
\usepackage{graphicx}
\graphicspath{{images/}}
\newcommand{\TallMath}[1]{$\displaystyle #1\rule[-1.4em]{0pt}{3.2em}$}
\newcommand{\CourseName}{Precalculus}
\newcommand{\MeetingLength}{55 minutes}
\newcommand{\UnitNumberName}{Unit 2: Polynomial Functions \quad}
\newcommand{\LessonNumberName}{Lesson 2.5: Polynomial Functions and End Behavior}

\begin{document}
\begin{center}
    {\Large\bfseries \CourseName \\
    {\small\bfseries \UnitNumberName \LessonNumberName}}
\end{center}

\begin{tcolorbox}[colback=sky, colframe=navy, boxrule=0.9pt, arc=2mm,
    left=2mm, right=2mm, top=1.5mm, bottom=1.5mm]
    \textbf{Primary Objective:} Students will describe the end behavior of polynomial functions
    using limit notation, and explain how the degree and sign of the leading coefficient determine
    what happens to output values as $x \to \pm\infty$.
    \hfill\textit{\small Skill: MP 3.A \quad (Big Idea: VAR)}
\end{tcolorbox}

\renewcommand{\arraystretch}{1.6}
\begin{skillbox}[Priority Ideas \& Skills]{goldbox}
\begin{tabularx}{\linewidth}{@{}X X@{}}
\textbf{AP Skills} &
\textbf{Key Understandings} \\
\midrule
\textbf{MP 3.A} --- Describe characteristics of a function with varying
levels of precision using appropriate mathematical language and notation.
(LO 1.6.A) &
\begin{itemize}[leftmargin=1.2em,itemsep=2pt,topsep=0pt]
  \item As $x \to \infty$ or $x \to -\infty$, output values of a nonconstant polynomial either
        increase or decrease without bound. (EK 1.6.A.1--2)
  \item End behavior is determined solely by the degree and sign of the leading coefficient. (EK 1.6.A.3)
  \item The leading term dominates all lower-degree terms for large $|x|$. (EK 1.6.A.3)
  \item Even-degree: both ends go the same direction; odd-degree: ends go opposite directions.
  \item Standard notation: $\lim_{x \to \infty} p(x) = \infty$ or $\lim_{x \to \infty} p(x) = -\infty$.
\end{itemize} \\
\end{tabularx}
\end{skillbox}

\begin{skillbox}[{Vocabulary, Concepts, \& Theorems}]{greenbox}
\renewcommand{\arraystretch}{1.4}
\begin{tabularx}{\linewidth}{@{} >{\leavevmode\bfseries\color{navy}}p{0.24\linewidth} X @{}}
End behavior       & What happens to $p(x)$ as $x$ increases or decreases without bound. \\
Leading term       & The term with the highest degree in a polynomial. \\
Leading coefficient & The coefficient of the leading term. \\
Limit notation     & $\lim_{x \to \infty} p(x) = \infty$ means output grows without bound as input increases. \\
Dominates          & The leading term dominates when its ratio to all other terms approaches 1 as $x \to \pm\infty$. \\
\end{tabularx}
\end{skillbox}

\begin{skillbox}[Activate Prior Knowledge \& Spiral Review]{sky}
    \begin{tabularx}{\linewidth}{@{}X c@{}}
        \begin{minipage}[t]{0.70\linewidth}\vspace{0pt}
            Please see attached warm-up (spiral review from Lessons 1.4--1.5).
            Skills reviewed: identifying degree and leading coefficient, writing a polynomial
            in factored form from given zeros and multiplicities, real vs.\ complex zeros.
        \end{minipage}
        &
        \begin{minipage}[t]{0.24\linewidth}\vspace{0pt}\centering
            \textit{\small (authored warm-up compiles to target/)}
        \end{minipage}
    \end{tabularx}
\end{skillbox}

\begin{skillbox}[Hook]{sky}
    \textbf{Rubber Band Ball:} A student keeps wrapping rubber bands around a ball. Suppose the
    radius (in mm) is modeled by $r(t) = 2t^3 - 50t^2 + 100t$, where $t$ is time in minutes.

    \medskip
    \textit{Discussion prompt:} As $t$ gets very large (say $t = 1{,}000$ minutes), does the
    $-50t^2$ term matter much compared to $2t^3$? At $t = 1{,}000$: $2t^3 = 2 \times 10^9$ but
    $-50t^2 = -5 \times 10^7$ --- less than $3\%$ of the cubic term. Which term ``wins'' for large $t$?

    \medskip
    Tell students: today we formalize exactly this idea --- the leading term always wins.
\end{skillbox}

\begin{skillbox}[Lesson]{sky}
\begin{multicols}{2}

\textbf{Part 1: Leading Term Dominance} (10 min)

Use $p(x) = x^3 - 100x^2$ to show the leading term dominates.

\smallskip
\renewcommand{\arraystretch}{1.3}
\begin{tabular}{|c|r|r|}
\hline
$x$ & $p(x)$ & $x^3$ \\
\hline
$10$    & $-9{,}000$ & $1{,}000$ \\
$100$   & $0$ & $1{,}000{,}000$ \\
$1{,}000$ & $9 \times 10^8$ & $10^9$ \\
\hline
\end{tabular}
\smallskip

Ask: At $x = 1{,}000$, the $-100x^2 = -10^8$ is large on its own, but $x^3 = 10^9$ is 10$\times$ bigger.
The ratio $-100x^2/x^3 = -100/x \to 0$.
\textbf{Conclusion:} Leading term $x^3$ controls end behavior.

\columnbreak

\textbf{Part 2: The Four End-Behavior Cases} (12 min)

\smallskip
\renewcommand{\arraystretch}{1.4}
\begin{tabular}{|c|c|c|c|}
\hline
\textbf{Degree} & \textbf{Lead $a_n$} & $x \to \infty$ & $x \to -\infty$ \\
\hline
Even & $>0$ & $\infty$  & $\infty$ \\
Even & $<0$ & $-\infty$ & $-\infty$ \\
Odd  & $>0$ & $\infty$  & $-\infty$ \\
Odd  & $<0$ & $-\infty$ & $\infty$ \\
\hline
\end{tabular}
\smallskip

Key insight: even-degree polynomials behave like $x^2$ (both ends same direction);
odd-degree behave like $x^3$ (ends go opposite directions).

\medskip
\textbf{Part 3: Using Limit Notation} (8 min)

Introduce $\lim_{x \to \infty}$ and $\lim_{x \to -\infty}$.

Example: $p(x) = -3x^4 + 2x^2 - 7$.
Leading term: $-3x^4$ (even, negative).

\[
  \lim_{x \to \infty} p(x) = -\infty \qquad \lim_{x \to -\infty} p(x) = -\infty
\]

Practice 2--3 examples chorally before the activity.

\end{multicols}
\end{skillbox}

\begin{skillbox}[Active Monitoring]{redbox}
While students work on guided notes and practice, circulate and check:
\begin{itemize}
  \item \textbf{Common error:} Students confuse ``the function has negative output values'' with
        ``end behavior is $-\infty$.'' Cold-call: ``Does a negative leading coefficient mean
        ALL outputs are negative?'' (No --- zeros and sign changes still exist.)
  \item \textbf{Common error:} Students use the constant or a middle term instead of the leading term.
        Ask: ``Which term matters when $x$ is a billion?''
  \item \textbf{Notation check:} Students should write $\lim_{x \to \infty} p(x) = \infty$ (not just
        ``goes to infinity''). AP exam requires limit notation.
  \item \textbf{Cold-call:} ``Give me a polynomial with $\lim_{x \to \infty} = -\infty$ and
        $\lim_{x \to -\infty} = \infty$.'' (Any odd-degree, negative-leading-coeff example.)
\end{itemize}
\end{skillbox}

\begin{skillbox}[Group Work \& Differentiation]{redbox}
Students work on the \textbf{Group Activity: End Behavior Explorer} in leveled tiers.
Each tier uses the same set of four polynomial functions but increases in cognitive demand.

\begin{itemize}
  \item \textbf{Tier R (Remediate):} Leading term provided for each polynomial; students use the
        rules table to fill in $\lim_{x \to \infty}$ and $\lim_{x \to -\infty}$ in limit notation.
  \item \textbf{Tier A (Apply):} Students identify the leading term from the full polynomial
        expression, predict end behavior, write limit notation, and verify with a calculator.
  \item \textbf{Tier E (Extend):} Students create their own polynomial with even degree, negative
        leading coefficient, and exactly 3 real zeros; then describe its end behavior fully.
\end{itemize}
\end{skillbox}

\begin{skillbox}[Individual Work \& Assessment]{redbox}
\textbf{Exit Ticket} (last 5--7 minutes, individual, collected):
\begin{enumerate}
  \item Given $p(x) = -2x^5 + 7x^3 - x + 3$, write $\lim_{x \to \infty} p(x)$ and $\lim_{x \to -\infty} p(x)$.
  \item A polynomial $q$ has degree 4 and positive leading coefficient. Complete: as $x \to \infty$,
        $q(x) \to \underline{\phantom{xxx}}$ and as $x \to -\infty$, $q(x) \to \underline{\phantom{xxx}}$.
\end{enumerate}
\textbf{Assessment note:} Look for correct identification of the leading term (odd degree, negative
coefficient in \#1) and proper limit notation. Students who write $-\infty$ and $\infty$ respectively
for \#1 demonstrate mastery of how odd-degree/negative leading coefficient determines opposite-direction ends.
\end{skillbox}

\begin{skillbox}[Reinforcement \& Extension]{goldbox}
\textbf{Homework:} 5 problems --- end behavior from leading term for 4 polynomials, verbal description
from a leading term, a conceptual question about parity, a table-based end-behavior inference, and one
AP-style multiple-choice item.

\textbf{Extension (optional):} Students investigate $p(x) = x^n$ for $n = 1, 2, 3, 4$ numerically
at $x = \pm100$ and generalize the four end-behavior patterns from the data.

\textbf{Preview:} Lesson 1.7 extends end behavior to \textit{rational} functions, where division by a
polynomial creates vertical asymptotes and horizontal (or oblique) asymptotes.
\end{skillbox}
\end{document}
